% =============================================================================
%                              CALCULUS PLUGIN
% =============================================================================
%
% Differential calculus, integral notation, derivative fractions, and vector
% calculus operators.
%
% Requires:
%   xamsmath core helpers. The braces and eqmicrotype options are optional
%   because this plugin provides fallbacks for their public helpers.

% =============================================================================
%                              FALLBACK SUPPORT
% =============================================================================

% -----------------------------------------------------------------------------
% User Commands
% -----------------------------------------------------------------------------
%
% \DeclareBracedMathOperator
%
% Purpose:
%   Provides ordinary AMS operator declarations when the braces module has not
%   already supplied the richer braced-operator implementation.
%
% Parameters:
%   *  - Optional star requests limit-style operator placement.
%   #1 - Operator command to define.
%   #2 - Operator text to display.
%
% Usage:
%   \DeclareBracedMathOperator{\grad}{grad}

\ProvideDocumentCommand{\DeclareBracedMathOperator}{smm}{
  \IfBooleanTF{#1}{\DeclareMathOperator*{#2}{#3}}{\DeclareMathOperator{#2}{#3}}
}

% -----------------------------------------------------------------------------
% User Commands
% -----------------------------------------------------------------------------
%
% \sfrac
%
% Purpose:
%   Provides a regular-fraction fallback when xfrac was not loaded by the
%   eqmicrotype option.
%
% Parameters:
%   #1 - Numerator.
%   #2 - Denominator.
%
% Usage:
%   \sfrac{dy}{dx}

\ProvideDocumentCommand{\sfrac}{mm}{\frac{#1}{#2}}

% =============================================================================
%                          COMMON CALCULUS NOTATION
% =============================================================================

% -----------------------------------------------------------------------------
% User Commands
% -----------------------------------------------------------------------------
%
% \interval
%
% Purpose:
%   Declares angle-bracket interval notation separated by a semicolon.
%
% Parameters:
%   #1 - Left endpoint.
%   #2 - Right endpoint.
%
% Usage:
%   \interval{a}{b}

\DeclarePairedDelimiterX\interval[2]{\langle}{\rangle}{#1;#2}

% -----------------------------------------------------------------------------
% User Commands
% -----------------------------------------------------------------------------
%
% \sign
%
% Purpose:
%   Typesets the signum operator.
%
% Usage:
%   \sign(x)

\DeclareMathOperator{\sign}{sign}

% -----------------------------------------------------------------------------
% User Commands
% -----------------------------------------------------------------------------
%
% \supp
%
% Purpose:
%   Typesets the support operator.
%
% Usage:
%   \supp(f)

\DeclareMathOperator{\supp}{supp}

% -----------------------------------------------------------------------------
% User Commands
% -----------------------------------------------------------------------------
%
% \dif
%
% Purpose:
%   Typesets a small upright differential symbol with integral spacing.
%
% Usage:
%   \int f(x) \dif x

\newmathcommand{\dif}{\mathop{}\!d}

% -----------------------------------------------------------------------------
% User Commands
% -----------------------------------------------------------------------------
%
% \Dif
%
% Purpose:
%   Typesets an uppercase differential operator for derivative notation.
%
% Usage:
%   \Dif f

\newmathcommand{\Dif}{\mathop{}\!D}

% =============================================================================
%                            DERIVATIVE TEMPLATE
% =============================================================================

\ExplSyntaxOn

% -----------------------------------------------------------------------------
% Internal Templates
% -----------------------------------------------------------------------------
%
% \tmpl_derivative_frac
%
% Purpose:
%   Renders PIE-aware derivative fractions with configurable differential symbol
%   and fraction command.
%
% Template parameters:
%   #1 - Differential symbol, such as \dif or \partial.
%   #2 - Fraction command, such as \frac, \sfrac, \tfrac, or \dfrac.
%
% Runtime parameters:
%   #3 - Prime tokens absorbed by mathcommand. Must be empty for derivatives.
%   #4 - Index tokens absorbed by mathcommand. Used as an evaluation suffix.
%   #5 - Exponent tokens absorbed by mathcommand. Used as derivative order.
%   #6 - Numerator expression.
%   #7 - Denominator variable or expression.
%
% Usage:
%   Internal template for \od, \pd, and their fraction-style variants.

\NewTemplateCommand\tmpl_derivative_frac<mm>{mmm mm}{
    \IfEmptyTF{#3}{
        \mleft.#2{#1#5#6}{#1#7#5}\IfEmptyTF{#4}{\mright.}{\mright|}#4
    }{
        \PackageError{xamsmath}{Primes~are~not~allowed~to~be~before~\verb|od|-like~commands}{}
    }
}

% =============================================================================
%                         ORDINARY DERIVATIVES
% =============================================================================

% -----------------------------------------------------------------------------
% User Commands
% -----------------------------------------------------------------------------
%
% \od
%
% Purpose:
%   Typesets an ordinary derivative using \frac.
%
% Parameters:
%   #1 - Numerator expression.
%   #2 - Denominator variable.
%
% PIE usage:
%   Superscripts specify derivative order. Subscripts specify evaluation points.
%
% Usage:
%   \od{y}{x}
%   \od^2_{x=0}{y}{x}

\TemplateInstancePIE\od\tmpl_derivative_frac<\dif, \frac>

% -----------------------------------------------------------------------------
% User Commands
% -----------------------------------------------------------------------------
%
% \lod
%
% Purpose:
%   Typesets an ordinary derivative using \sfrac.
%
% Parameters:
%   #1 - Numerator expression.
%   #2 - Denominator variable.
%
% Usage:
%   \lod{y}{x}

\TemplateInstancePIE\lod\tmpl_derivative_frac<\dif, \sfrac>

% -----------------------------------------------------------------------------
% User Commands
% -----------------------------------------------------------------------------
%
% \tod
%
% Purpose:
%   Typesets an ordinary derivative using \tfrac.
%
% Parameters:
%   #1 - Numerator expression.
%   #2 - Denominator variable.
%
% Usage:
%   \tod{y}{x}

\TemplateInstancePIE\tod\tmpl_derivative_frac<\dif, \tfrac>

% -----------------------------------------------------------------------------
% User Commands
% -----------------------------------------------------------------------------
%
% \dod
%
% Purpose:
%   Typesets an ordinary derivative using \dfrac.
%
% Parameters:
%   #1 - Numerator expression.
%   #2 - Denominator variable.
%
% Usage:
%   \dod{y}{x}

\TemplateInstancePIE\dod\tmpl_derivative_frac<\dif, \dfrac>

% =============================================================================
%                          PARTIAL DERIVATIVES
% =============================================================================

% -----------------------------------------------------------------------------
% User Commands
% -----------------------------------------------------------------------------
%
% \pd
%
% Purpose:
%   Typesets a partial derivative using \frac.
%
% Parameters:
%   #1 - Numerator expression.
%   #2 - Denominator variable or mixed-partial expression.
%
% Usage:
%   \pd{u}{t}
%   \pd^2{u}{x \partial y}

\TemplateInstancePIE\pd\tmpl_derivative_frac<\partial, \frac>

% -----------------------------------------------------------------------------
% User Commands
% -----------------------------------------------------------------------------
%
% \lpd
%
% Purpose:
%   Typesets a partial derivative using \sfrac.
%
% Parameters:
%   #1 - Numerator expression.
%   #2 - Denominator variable or mixed-partial expression.
%
% Usage:
%   \lpd{u}{t}

\TemplateInstancePIE\lpd\tmpl_derivative_frac<\partial, \sfrac>

% -----------------------------------------------------------------------------
% User Commands
% -----------------------------------------------------------------------------
%
% \tpd
%
% Purpose:
%   Typesets a partial derivative using \tfrac.
%
% Parameters:
%   #1 - Numerator expression.
%   #2 - Denominator variable or mixed-partial expression.
%
% Usage:
%   \tpd{u}{t}

\TemplateInstancePIE\tpd\tmpl_derivative_frac<\partial, \tfrac>

% -----------------------------------------------------------------------------
% User Commands
% -----------------------------------------------------------------------------
%
% \dpd
%
% Purpose:
%   Typesets a partial derivative using \dfrac.
%
% Parameters:
%   #1 - Numerator expression.
%   #2 - Denominator variable or mixed-partial expression.
%
% Usage:
%   \dpd{u}{t}

\TemplateInstancePIE\dpd\tmpl_derivative_frac<\partial, \dfrac>

\ExplSyntaxOff

% =============================================================================
%                       MULTIVARIABLE CALCULUS OPERATORS
% =============================================================================

% -----------------------------------------------------------------------------
% User Commands
% -----------------------------------------------------------------------------
%
% \grad
%
% Purpose:
%   Typesets the gradient operator.
%
% Usage:
%   \grad f
%   \grad(f)

\DeclareBracedMathOperator{\grad}{grad}

% -----------------------------------------------------------------------------
% User Commands
% -----------------------------------------------------------------------------
%
% \rot
%
% Purpose:
%   Typesets the rotation or curl operator.
%
% Usage:
%   \rot \vec{F}

\DeclareBracedMathOperator{\rot}{rot}

% -----------------------------------------------------------------------------
% User Commands
% -----------------------------------------------------------------------------
%
% \div
%
% Purpose:
%   Typesets the divergence operator, replacing TeX's division symbol command.
%
% Usage:
%   \div \vec{F}

\let\div\relax
\DeclareBracedMathOperator{\div}{div}

% =============================================================================
%                                USAGE NOTES
% =============================================================================
%
% Common notation:
%   \interval{a}{b}
%   \sign(x)
%   \supp(f)
%   \int f(x) \dif x
%
% Ordinary derivatives:
%   \od{y}{x}
%   \od^2_{x=0}{y}{x}
%   \lod{y}{x}
%
% Partial derivatives:
%   \pd{u}{t}
%   \pd^2{u}{x \partial y}
%   \tpd{u}{t}
%
% Multivariable calculus:
%   \grad f
%   \rot \vec{F}
%   \div \vec{F}
